% !TeX root = ../Book2.tex
\BookPartGuide[附录导读]

附录从正文中几个值得继续追问的问题展开：径向对称怎样改变 Fourier 公式，特殊函数怎样成为正交基，粗糙系数下怎样比较正解，以及频率方法如何进入非线性演化。可以从正在研究的问题选择相应专题，下面说明各附录的内容与阅读联系。

\ProseParagraphBreak
\noindent\textbf{附录 A：径向 Fourier 与 Bessel 工具。}
\cref{b2:thm:A2-sphere-fourier}给出球面测度的 Fourier 公式，\cref{b2:thm:A2-bessel-asymptotic}给出固定实阶的大参数渐近，供\chapref{b2:ch:23}研究球面衰减。\cref{b2:thm:A2-bochner-riesz}给出球内幂权的变换，\cref{b2:prop:A2-kernel-integrability}辨认可积核的准确阈值，供\chapref{b2:ch:23}和\chapref{b2:ch:24}回查。\chapref{b2:ch:23}将继续研究复参数变化时的统一增长控制。

\ProseParagraphBreak
\noindent\textbf{附录 B：特殊函数的谱模型。}
\chapref{b2:ch:05}之后可读 Hermite，区间谱理论之后可读 Legendre 与球谐。进入 Hermite 部分前，若曾略过选读的\subsecref{b2:subsec:050204}，须先回读其正交完备性证明；本附录核对归一化和 Rodrigues 表达，并研究谐振子的自伴定义域、核与紧逆。Legendre 和球谐的完备性在各自模型中建立；形式特征方程本身不保证可以展开一般数据。

\ProseParagraphBreak
\noindent\textbf{附录 C：主值、原子空间与加权理论。}
\chapref{b2:ch:17}和\chapref{b2:ch:18}之后可研究极大截断的振荡传递、John--Nirenberg 指数尾、原子 \(H^1\)、BMO 对偶和权。这里的原子 Hardy 空间 \(H^1\) 与 PDE 的一阶平方 Sobolev 空间须按定义区分。极大截断在 \(L^p\)、\(p>1\) 中的控制，使光滑输入上的主值收敛能够传到一般输入。研究 Harnack 比较时，可重点阅读\secref{b2:sec:app-C-2}：John--Nirenberg 估计把平均振荡控制转为指数可积性。

\ProseParagraphBreak
\noindent\textbf{附录 D：半群与生成元。}
接\chapref{b2:ch:12}热方程，把核卷积推广为抽象演化，区分温和解、强解和生成元定义域。一般非负自伴算子的谱积分采用另行注明的无界谱定理与函数演算，其范围大于\chapref{b2:ch:10}的紧自伴离散谱；读者可以采用这个一般谱定理，也可先从前面的具体半群模型入手。

\ProseParagraphBreak
\noindent\textbf{附录 E：色散估计与非线性模型。}
需要插值、势估计、驻相及所用函数空间。从 Schrödinger 色散核进入非端点 Strichartz 估计，再处理一维三次 NLS 的固定点、守恒与延拓。时空混合范数把空间衰减与时间积累联系起来，为控制非线性项提供合适的解空间。

\ProseParagraphBreak
\noindent\textbf{附录 F：集中紧性与变分。}
接\chapref{b2:ch:14}和第一卷的 Sobolev 嵌入、弱紧性，比较质量消失、分裂与集中，并在列明的次临界模型中构造极小元。它说明直接法失去强紧性后还需要控制什么。

\ProseParagraphBreak
\noindent\textbf{附录 G：相位与算子的微局部分析。}
以\chapref{b2:ch:05}分布、\chapref{b2:ch:21}振荡积分及相关算子估计为基础。波前集同时记录奇性的位置与频率方向；符号演算和椭圆参数逆刻画正则性，Fourier 积分算子的相位则描述奇性怎样传播。

\ProseParagraphBreak
\noindent\textbf{附录 H：Harnack 证明与非线性能量。}
\secref{b2:sec:H201}接\chapref{b2:ch:16}的实标量散度型方程。普通主线可先接受\cref{b2:thm:16-measurable-harnack}，继续其 Hölder 推论；补齐证明时，先读\secref{b2:sec:1701}中二进立方体与分化所需部分，再读\cref{b2:thm:C-john-nirenberg}，随后读幂迭代、对数估计及\cref{b2:thm:H2-linear-harnack}。\secref{b2:sec:H202}的 \(p\)-Laplace 变分更接近\chapref{b2:ch:14}，非线性增长将决定新的截断估计与局部上界。

\ProseParagraphBreak
\noindent\textbf{附录 I：平移、对偶与群上的 Fourier 分析。}
周期、离散和连续频率为什么有不同的形状？\appref{b2:app:I}从 Haar 积分与字符回答这个问题，把本卷的 Fourier 级数、卷积和 Plancherel 理论放到共同的群结构中。基本模型逐项计算测度与反演常数，仿射群说明左右平移的差异，紧群的矩阵系数则展示非交换频率。已熟悉测度积分及本卷初等 Fourier 理论的读者，可从这些模型进入一般对偶定理。
